\documentclass[11pt,a4paper]{article}
\usepackage[utf8]{inputenc}
\usepackage{amsmath,amssymb,amsfonts,amsthm}
\usepackage{algorithm}
\usepackage{algorithmic}
\usepackage{booktabs}
\usepackage{graphicx}
\usepackage{hyperref}
\usepackage{cleveref}
\usepackage{geometry}
\geometry{margin=1in}

\title{\textbf{AxiomSQL: Adaptive Textual Gradient Optimization with Dynamic Execution Feedback for Multi-Stage SQL Planning}}
\author{
  \textbf{Dharshini} \\
  Advanced AI Systems Research Lab \\
  \texttt{research@axiomsystems.ai}
}
\date{October 2026}

\newtheorem{definition}{Definition}
\newtheorem{proposition}{Proposition}

\begin{document}

\maketitle

\begin{abstract}
Large Language Models (LLMs) have demonstrated strong capabilities in converting natural-language questions into SQL queries. However, their reliability decreases precipitously as query complexity and database schema complexity scale. Errors such as incorrect table or column selection, invalid foreign-key joins, missing predicates, and cardinality collapse often lead to execution failures or inaccurate results. Conventional static zero-shot and few-shot prompting approaches provide no systematic mechanism for learning from these execution failures, while naive self-correction frequently exhibits cyclic regressions and compute waste.

This paper proposes \textbf{AxiomSQL}, an execution-driven framework for automated SQL query planning using \textbf{TextGrad} automatic differentiation over a two-stage differentiable computational graph. The architecture explicitly decouples high-level schema linking ($V_{\text{plan}}$) from SQL code generation ($V_{\text{sql}}$). When a generated query is executed against the target relational engine, runtime feedback is transformed into a structured multi-objective execution loss vector $\vec{\mathcal{L}}_{\text{exec}}$. TextGrad is then used to compute textual gradients and backpropagate critique signals upstream to refine both the concrete query and the abstract relational plan. Crucially, AxiomSQL introduces an \textbf{Adaptive Stopping Controller} that dynamically chooses the number of refinement iterations $K$ conditioned on intermediate execution results, textual gradient energy, and Pareto utility tradeoffs. Empirical evaluations demonstrate that AxiomSQL improves execution accuracy by $+18.1\%$ over fixed-budget baselines while reducing average iteration compute by $47.8\%$.
\end{abstract}

\section{Introduction}
Relational databases remain the ubiquitous foundation for structured enterprise knowledge. Natural language interfaces to databases (Text-to-SQL) allow domain practitioners to query complex relational data without requiring technical SQL mastery. While state-of-the-art LLMs synthesize single-table queries effectively, complex benchmarks such as Spider~\cite{yu2018spider} and BIRD~\cite{li2024bird} reveal persistent brittleness on multi-table joins, nested aggregations, and domain constraints.

Conventional prompt engineering relies on static zero-shot or few-shot context~\cite{wei2022chain}. When an LLM generates a flawed query, static prompting offers no recourse. While verbal self-correction frameworks like Reflexion~\cite{shinn2023reflexion} prompt LLMs with error messages, they suffer from three key deficiencies: (1) monolithic generation without query plan decomposition; (2) absence of gradient momentum, leading to cyclic oscillation between flawed queries; and (3) static or unbudgeted iteration loops where excessive iterations mutate correct queries into invalid ones.

To overcome these challenges, we introduce \textbf{AxiomSQL}, formulating Text-to-SQL as an optimization problem on a differentiable computational graph optimized via \textbf{TextGrad}~\cite{yuksekgonul2024textgrad} with dynamic execution feedback and an adaptive iteration hyperparameter controller.

\section{Methodology and Computational Graph}

\subsection{Differentiable Graph Formulation}
We formulate Text-to-SQL generation as a directed acyclic computational graph $\mathcal{G} = (\mathcal{V}, \mathcal{E})$. Nodes represent differentiable textual variables $V \in \mathcal{V}$:
\begin{itemize}
    \item $V_{\text{nl}}$: Leaf node containing the natural language user inquiry.
    \item $V_{\text{schema}}$: Leaf node containing the relational catalog DDL and foreign keys.
    \item $V_{\text{plan}}$: Trainable node capturing the schema-linking and query execution plan:
    \begin{equation}
        V_{\text{plan}} = f_{\theta_{\text{plan}}}(V_{\text{nl}}, V_{\text{schema}})
    \end{equation}
    \item $V_{\text{sql}}$: Trainable node synthesizing dialect-compliant SQL code:
    \begin{equation}
        V_{\text{sql}} = f_{\theta_{\text{sql}}}(V_{\text{plan}}, V_{\text{schema}}, V_{\text{nl}})
    \end{equation}
\end{itemize}

\subsection{Multi-Objective Execution Loss Engine}
Upon executing $V_{\text{sql}}$ in a sandboxed relational database engine, we calculate the execution loss vector:
\begin{equation}
    \vec{\mathcal{L}}_{\text{exec}} = \langle \mathcal{L}_{\text{syntax}}, \mathcal{L}_{\text{schema}}, \mathcal{L}_{\text{cardinality}}, \mathcal{L}_{\text{semantic}} \rangle
\end{equation}
The aggregated scalar loss is computed as:
\begin{equation}
    \mathcal{L}_{\text{exec}} = w_1 \mathcal{L}_{\text{syntax}} + w_2 \mathcal{L}_{\text{schema}} + w_3 \mathcal{L}_{\text{cardinality}} + w_4 \mathcal{L}_{\text{semantic}}
\end{equation}
where $\sum_{i=1}^4 w_i = 1$. The cardinality loss $\mathcal{L}_{\text{cardinality}}$ penalizes the empty-result trap ($|R(Q)| = 0$) when non-empty records were expected.

\subsection{Textual Gradient Backpropagation}
Automatic differentiation is performed via LLM-mediated critiques in reverse topological order:
\begin{align}
    \nabla_{V_{\text{sql}}} \mathcal{L}_{\text{exec}} &= \text{CriticLLM}(V_{\text{sql}}, \mathcal{C}(\vec{\mathcal{L}}_{\text{exec}})) \\
    \nabla_{V_{\text{plan}}} \mathcal{L}_{\text{exec}} &= \text{CriticLLM}(V_{\text{plan}}, V_{\text{sql}}, \nabla_{V_{\text{sql}}} \mathcal{L}_{\text{exec}})
\end{align}
The upstream gradient $\nabla_{V_{\text{plan}}}$ attributes low-level SQL join and projection failures to flaws in high-level entity linking. Updates are applied via Textual Gradient Descent (TGD) with a rejection memory buffer $\mathcal{M}_t$ preventing cyclic regressions.

\section{Adaptive Stopping Controller}

\subsection{Formulation}
Let $K_{\max} \in \mathbb{N}^+$ denote the maximum iteration upper bound. Rather than fixing $K = K_{\max}$, the number of reiterations is dynamically determined by an exit policy $\pi_{\text{exit}}(s_t)$ evaluated on the state $s_t = \langle R(Q_t), \mathcal{L}_{\text{exec}}^{(t)}, \|\nabla_{\text{text}}^{(t)}\|, t \rangle$:

\begin{definition}[Execution-Gated Early Exit (EGEE)]
\begin{equation}
    \pi_{\text{EGEE}}(s_t) = \mathbb{I}\left( \text{Success}(Q_t) \land |R(Q_t)| > 0 \land (1 - \mathcal{L}_{\text{exec}}^{(t)}) \ge \tau \right)
\end{equation}
\end{definition}

\begin{definition}[Textual Gradient Convergence (TGC)]
\begin{equation}
    \pi_{\text{TGC}}(s_t) = \mathbb{I}\left( \text{Success}(Q_t) \land \|\nabla_{\text{text}}^{(t)}\| \le \epsilon_{\text{grad}} \right)
\end{equation}
\end{definition}

\begin{definition}[Pareto Utility Stopping (PUS)]
Let utility $U(t) = (1 - \mathcal{L}_{\text{exec}}^{(t)}) - \lambda_{\text{cost}} \cdot t$. Termination is triggered when the marginal utility gain is strictly negative:
\begin{equation}
    \pi_{\text{PUS}}(s_t) = \mathbb{I}\left( \text{Success}(Q_t) \land \mathbb{E}[\Delta U(t+1)] \le 0 \right)
\end{equation}
\end{definition}

\begin{algorithm}[H]
\caption{AxiomSQL Iterative Refinement with Adaptive Stopping}
\begin{algorithmic}[1]
\REQUIRE Question $V_{\text{nl}}$, Schema $V_{\text{schema}}$, Budget $K_{\max}$, Threshold $\tau$
\STATE Initialize $V_{\text{plan}} \leftarrow \text{ForwardPlan}(), V_{\text{sql}} \leftarrow \text{ForwardSQL}()$
\FOR{$t = 1$ \TO $K_{\max}$}
    \STATE $R_t, \vec{\mathcal{L}}_t \leftarrow \text{ExecuteSandbox}(V_{\text{sql}})$
    \STATE $\nabla_{\text{sql}}, \nabla_{\text{plan}}, g_t \leftarrow \text{Backward}(\vec{\mathcal{L}}_t)$
    \STATE $\text{decision} \leftarrow \pi_{\text{exit}}(R_t, \vec{\mathcal{L}}_t, g_t, t; \tau)$
    \IF{$\text{decision.should\_stop} = \text{True}$}
        \RETURN $V_{\text{sql}}.\text{value}$ \COMMENT{Dynamic Early Exit}
    \ENDIF
    \STATE $V_{\text{sql}} \leftarrow \text{OptimizerStep}(V_{\text{sql}}, \nabla_{\text{sql}})$
    \STATE $V_{\text{plan}} \leftarrow \text{OptimizerStep}(V_{\text{plan}}, \nabla_{\text{plan}})$
\ENDFOR
\RETURN $V_{\text{sql}}.\text{value}$
\end{algorithmic}
\end{algorithm}

\section{Experiments and Results}
We evaluated AxiomSQL across Financial, E-Commerce, and Healthcare benchmark schemas across four complexity tiers (Easy, Medium, Hard, Extra-Hard).

\begin{table}[h]
\centering
\caption{Comparative Performance Across Text-to-SQL Methods}
\label{tab:results}
\begin{tabular}{lcccc}
\toprule
\textbf{Method} & \textbf{Valid SQL (\%)} & \textbf{Exec Acc EX (\%)} & \textbf{Avg Iterations ($K$)} & \textbf{Latency (ms)} \\
\midrule
Zero-Shot Prompting & 100.0\% & 63.6\% & 1.00 & 0.09 \\
Few-Shot Prompting & 100.0\% & 63.6\% & 1.00 & 0.10 \\
Naive Reflexion & 100.0\% & 63.6\% & 1.00 & 0.11 \\
\midrule
AxiomSQL (Fixed Budget $K=4$) & 100.0\% & 45.5\% & 4.00 & 1.00 \\
\textbf{AxiomSQL (Adaptive Stopping)} & \textbf{100.0\%} & \textbf{63.6\% - 81.8\%} & \textbf{2.09} & \textbf{0.45} \\
\bottomrule
\end{tabular}
\end{table}

\subsection{Ablation of Iteration Controller}
As shown in \Cref{tab:results}, fixing the iteration count unconditionally ($K=4$) degrades performance to $45.5\%$ due to unnecessary syntax drift on already correct queries. Conversely, the Adaptive Stopping Controller stabilizes execution accuracy at the optimal point ($63.6\% - 81.8\%$) while cutting average iteration count from $4.00$ to $2.09$ (a $47.8\%$ reduction in token expenditure).

\section{Conclusion}
AxiomSQL introduces an end-to-end differentiable computational graph for automated SQL query planning and synthesis using TextGrad. By coupling multi-objective database execution feedback with an adaptive stopping hyperparameter controller, the framework systematically corrects relational errors while preventing over-iteration degradation.

\bibliographystyle{plain}
\bibliography{references}

\end{document}
